% !TEX root = ../b-rg.tex
\chapter{Relative Clauses}\label{ch:relative-clauses}

Borlish has no relative pronoun for subjects or objects. A relative clause on a subject or an object is a participial phrase placed after the head noun: a present participle for a subject, and a past participle, usually with an agentive pronoun, for an object. Finite relative clauses are used only for possessors and partitives, with \bl{dont}, and for places, with \bl{yon} `where'. In informal speech \bl{yon} has also spread to other relations. The participial relative is by far the commonest type, in speech and writing alike.

The strategies are summarised below and described in the sections that follow.
\begin{itemize}
    \item \textbf{Subject}: head + present participle (\bl{y domn hamtant cas coster} `the woman who lives next door'), or a perfect participle for an earlier event (Section~\ref{sec:subject-relatives}).
    \item \textbf{Object}: head + agentive pronoun + past participle (\bl{toð gent meyon connuð} `everyone I know'), or head + past participle + \bl{de}/\bl{par} + noun phrase (Section~\ref{sec:object-relatives}).
    \item \textbf{Oblique}: as for objects, with the preposition left at the end in its adverbial form (\bl{y caðer teyon asseïð dessur} `the chair you're sitting on'; Section~\ref{sec:oblique-relatives}).
    \item \textbf{Place}: \bl{yon} + finite clause (\bl{y majon yon jo fo nað} `the house where I was born'; Section~\ref{sec:finite-relatives}).
    \item \textbf{Possessor and partitive}: \bl{dont} + finite clause (\bl{y segr dont y can me mordoy} `the man whose dog bit me').
    \item \textbf{Free relatives}: \bl{l'oç} with a participial clause, and \bl{cuscon}, \bl{calscon} (Section~\ref{sec:free-relatives}).
\end{itemize}

\section{Subject relatives}\label{sec:subject-relatives}
A relative clause on a subject is a present participle following the head noun, with its own objects and adjuncts. There is no relative word and no complementiser.
    \glsen{Y domn hamtant cas coster es sanat.}{i ˈdɔ.mn̩ hamˈtant kaz kɔˈstɛr ɛz saˈnat}{{\Def} woman {dwell-\Prs.\Ptcp} at side be doctor}{The woman who lives next door is a doctor.}
The participle carries no tense. A past situation is placed in the past by the context or by a time adverb such as \bl{lor} `then' or \bl{paravant} `earlier'. For a process in progress at the time, the periphrasis \bl{a mesur de} `in the course of' may be used: \bl{y segr a mesur lor de parolar a Halina} `the man who was in the middle of talking to Halina then'.
    \glsen{Y segr lor parolant a Halina es mell'yam.}{i ˈsij.r̩ lɔr ˌpa.roˈlant a ˌha.liˈna ɛz mɛˈljam}{{\Def} man then {speak-\Prs.\Ptcp} to Halina be {\First\Sg.\Poss-\Def=uncle}}{The man who was talking to Halina is my uncle.}
An event that precedes the time of the main clause is expressed with the perfect participle: \bl{ayent} with the past participle for a transitive verb, \bl{stant} for an intransitive one (Section~\ref{sec:perfect-star}). As in the finite perfect, the object stands between the two parts: \bl{l'un ayent y fenestr fragnt} `the one who had broken the window', \bl{y segr stant parolað a Halina} `the man who had spoken to Halina'.
    \glsen{Jovanna sta l'un ayent y fenestr fragnt.}{ˌʒo.vaˈna sta lɪn eˈjɛnt i feˈnɛs.tr̩ frɛjnt}{Jovanna {be.\Ipfv} {\Def=one} {have-\Prs.\Ptcp} {\Def} window {break.\Pst.\Ptcp}}{It was Jovanna who broke the window.}
This last example is a cleft, in which the relative clause modifies \bl{l'un} `the one' (Chapter~\ref{ch:information-structure}).
% CHECK: is a finite clause with ig ever possible for a subject relative (*y
% domn ig hamt cas coster)? No example has been found or elicited.

\section{Object relatives}\label{sec:object-relatives}

\subsection{With an agentive pronoun}
A relative clause on a direct object consists of the agent and a past participle. When the agent is a pronoun, it takes the agentive form (\bl{meyon, teyon, noscon, voscon, seyon}; Section~\ref{sec:agentive}), which stands immediately before the participle. \bl{Seyon} serves the third person in both numbers, and \bl{voscon} is used for the polite singular as well as the plural.
    \glsen{La cant ottem a toð gent meyon connuð.}{la kant ˈɔ.tɛm a tɔθ ʒɛnt miˈjɔn kɔˈnɪθ}{{\Third\Sg.\F} sing better than all people {\First\Sg.\Agt} {know.\Pst.\Ptcp}}{She sings better than anyone I know.}
    \glsen{My fraðr dessoclau y castel blader meyon costroit.}{mi ˈfra.ðr̩ ˌde.soˈklo i kaˈstɛl blaˈdɛr miˈjɔn koˈstrɔjt}{{\First\Sg.\Poss} brother {topple-\Pst} {\Def} castle {card-\Adj} {\First\Sg.\Agt} {build.\Pst.\Ptcp}}{My brother knocked over the house of cards I built.}
The participle is invariable and tenseless, like that of a subject relative. The time of the relative clause is understood from the context: \bl{y pog stoir seyon oïð} `the few tales he's heard', \bl{l'oç seyon sentið} `what he feels'.

Any other object of the participle, including a clitic, comes before it, after the agentive pronoun: \bl{y sovragl seyon me dað} `the belt she gave me'. A participial relative is negated with \bl{noc} before the agentive pronoun: \bl{l'oç noc meyon amað} `the one I don't like'.
    \glsen{Y soup es l'oç noc meyon amað.}{i sup ɛz lɔts nɔk miˈjɔn aˈmaθ}{{\Def} soup be {\Def=\Dist.\Sg} {\Neg} {\First\Sg.\Agt} {love-\Pst.\Ptcp}}{It's the soup that I don't like.}

\subsection{With a noun-phrase agent}
The agentive pronouns have no noun counterpart. When the agent is a full noun phrase, it follows the participle in a prepositional phrase with \bl{de} `of, from' or \bl{par} `by'. Speakers feel a slight difference between them. \bl{De} presents the agent as the author or source of what is described, and \bl{par} as the one who carried out the action: \bl{y lettr m'accaçað de my fraðr} suggests that the brother wrote the letter, and \bl{par my fraðr} that he posted it. In most contexts, however, the two are interchangeable.
    \glsen{Y lettr m'accaçað de my fraðr me parvenoy jamay.}{i ˈlɛ.tr̩ ˌma.kaˈtsaθ de mi ˈfra.ðr̩ me ˌpar.veˈnɔj ʒaˈme}{{\Def} letter {\First\Sg.\Obl=send-\Pst.\Ptcp} of {\First\Sg.\Poss} brother {\First\Sg.\Obl} {reach-\Pst} never}{The letter my brother sent me never arrived.}

\subsection{With the potential}
In place of the past participle, a relative clause may contain the potential in \bl{-abr} or \bl{-br}, `that can be \ldots{}ed', with the same agentive pronouns (Chapter~\ref{ch:word-formation}). It can be negated with the prefix \bl{in-}, as in \bl{un dragoman voscon inconfeyabr} `an interpreter you don't trust'.
    \glsen{Nos nolau hur dar for an svattour noscon confiabr.}{nɔz noˈlo hɪr dar fɔr an svaˈtur noˈxɔn kɔnˈfja.br̩}{{\First\Pl} {\Neg.\Fut-\Ipfv.\First\Pl} hire {give.\Inf} only {to.\Indf} detective {\First\Pl.\Agt} {trust-\Pot}}{We were only going to hire a PI we could trust.}
% CHECK: the agentive pronoun may be preceded by an optional oc in some
% relatives (y coutel (oc) meyon y pan trencað con, E134; y segr (oc) meyon
% parolað ig, E151). Is this oc the demonstrative, a relative marker, or a
% register variant? Which relatives allow it?

\section{Oblique relatives and stranding}\label{sec:oblique-relatives}
A relative clause on the object of a preposition is built in the same way as an object relative, with an agentive pronoun or a noun-phrase agent and a past participle. The preposition is not moved to the front of the clause. Instead it stays behind at the end of the clause, in the form of the corresponding adverb.
\begin{table}[ht]
    \centering
    \begin{tabular}{l l l l}
        \multicolumn{4}{c}{Prepositions and their stranded forms} \\
        \hline
        preposition & & stranded & \\
        \hline
        \bl{a} & `to, at' & \bl{ig} & `there, to it' \\
        \bl{sur} & `on' & \bl{dessur} & `on top' \\
        \bl{sou} & `under' & \bl{dessou} & `underneath' \\
        \bl{vars} & `towards' & \bl{avars} & `in that direction' \\
        \bl{con} & `with' & \bl{con} & `with it' \\
        \hline
    \end{tabular}
\end{table}
Only \bl{con} has the same form in both uses.
    \glsen{Y segr meyon parolað ig sta Fransc.}{i sijr̩ miˈjɔn ˌpa.roˈlaθ aj sta franx}{{\Def} man {\First\Sg.\Agt} {speak-\Pst.\Ptcp} {to.it} {be.\Ipfv} French}{The man I spoke to was French.}
    \glsen{Y caðer teyon asseïð dessur es ront.}{i kaˈðɛr tiˈjɔn ˌa.siˈjɪθ dɛˈsɪr ɛz rɔnt}{{\Def} chair {\Second\Sg.\Agt} {sit-\Pst.\Ptcp} {on.it} be broken}{The chair you're sitting on is broken.}
    \glsen{Yon es y coutel meyon y pan trencað con?}{jɔn ɛz i kuˈtɛl miˈjɔn i pan trɛnˈkaθ kɔn}{where be {\Def} knife {\First\Sg.\Agt} {\Def} bread {cut-\Pst.\Ptcp} {with.it}}{Where's the knife I cut the bread with?}
The participle's own direct object comes before it in the usual way (\bl{meyon y pan trencað con}), and the stranded adverb closes the clause. This is also the case where the verb would take a preposition with a finite verb: \bl{pesar a} `think about' gives \bl{y majon meyon pesað ig} `the house I was thinking about'. Stranding is confined to relative clauses. In a question the preposition is fronted with the question word (\bl{A cal es tu pesant?} `What are you thinking about?'; Chapter~\ref{ch:sentence-types}).
% CHECK: stranded forms of the other prepositions (de, ne, par, por, pre,
% des)? Is stranding possible with all of them?

\section{Finite relatives}\label{sec:finite-relatives}

\subsection{\bl{yon} `where'}
A relative clause on a place is a finite clause introduced by \bl{yon} `where', with the subject in its normal position before the verb, as in other embedded wh-clauses.
    \glsen{L'oc es y majon yon jo fo nað.}{lɔk ɛz i maˈʒɔn jɔn ʒo fo naθ}{{\Def=\Prox.\Sg} be {\Def} house where {\First\Sg} {be.\Pst} {bear.\Pst.\Ptcp}}{This is the house where I was born.}
In informal speech, \bl{yon} is also used as a general relative word with a finite clause, in place of a participial relative. A preposition is stranded in the same way as in a participial relative. Unlike the participle, the finite verb expresses tense and aspect, which can make this construction the more precise of the two in speech. It is not found in formal writing.
    \glsen{Y majon yon j'er pesant ig a se caut vençað.}{i maˈʒɔn jɔn ʒɛr peˈzant aj a se kot vɛnˈtsaθ}{{\Def} house where {\First\Sg=\Cop.\Ipfv} {think-\Prs.\Ptcp} {to.it} have {\Refl} {get.\Pst.\Ptcp} {sell-\Pst.\Ptcp}}{The house I was thinking about has been sold.}
Compare \bl{y caðer yon tu's t'asseyent dessur} `the chair you're sitting on' and \bl{y coutel yon jo trencau y pan con} `the knife I cut the bread with'.

\subsection{\bl{dont} `of which, whose'}
The relative pronoun \bl{dont} introduces finite relative clauses on a possessor or on the whole of which a part is singled out. In a possessive relative, \bl{dont} is followed immediately by the possessed noun with its article. Where that noun is the subject, the clause continues in the normal order.
    \glsen{L'oç es y segr dont y can me mordoy.}{lɔts ɛz i sijr dɔnt i kan me mɔrˈdɔj}{{\Def=\Dist.\Sg} be {\Def} man whose {\Def} dog {\First\Sg.\Obl} {bite-\Pst}}{That's the man whose dog bit me.}
Where the possessed noun is an object, it still comes first, and the subject and verb invert after it, as after any fronted constituent: \bl{un loy dont y lengaç n'ay jo nien} `a place of whose languages I know nothing'. Other examples are \bl{dont y fraðr Lohis fo mort} `whose brother Lohis had died' and \bl{luy es l'un dont y dors caf appoyað} `he is the one with his back up against the wall' (literally `whose back gets leant').

In a partitive relative, \bl{dont} is followed by a quantified noun phrase or a name, which picks out one or more members of the set: \bl{ogt ivan, dont Adagle sta treyem} `eight children, of whom Adaillé was the third'.
    \glsen{L'a dou fraðr, dont y dou son sanat.}{la du ˈfra.ðr̩ dɔnt i du sɔn saˈnat}{{\Third\Sg.\F=have} two brother {of.whom} {\Def} two {be.\Third\Pl} doctor}{She has two brothers, both of whom are doctors.}

\section{Non-restrictive relatives}\label{sec:nonrestrictive}
A non-restrictive relative clause is a participial phrase in apposition, set off by commas, with no relative word. With a predicate noun or adjective the participle is \bl{stant} `being'.
    \glsen{My vecin, stant sanat, es eut a Spagn.}{mi veˈtsɪn stant saˈnat ɛz awt a spɛjn}{{\First\Sg.\Poss} neighbour {be-\Prs.\Ptcp} doctor be {go.\Pst.\Ptcp} to Spain}{My neighbour, who is a doctor, has gone to Spain.}
The perfect participle is common in this use, and may be coordinated with a finite \bl{dont} clause: \bl{Raimond Aurellex, ayent sy sour y jundressem mariðað e dont y fraðr Lohis fo mort\ldots} `Raimond Aurellex, who had married his youngest sister and whose brother Lohis had died\ldots'.

\section{Free relatives}\label{sec:free-relatives}
A free relative `what, that which' is formed with the demonstrative pronoun \bl{l'oç} and a participial relative (Section~\ref{sec:dem-without-art}): \bl{l'oç seyon rovað} `what he asked', \bl{l'oç seyon creut} `what he believed'. A person is \bl{l'un} `the one' in place of \bl{l'oç} (\bl{l'un meyon vis ag quarf} `the one I saw at the station'). Both are the basis of the cleft and pseudocleft constructions (Chapter~\ref{ch:information-structure}).
    \glsen{L'oç meyon recas es un cogmað lonc.}{lɔts miˈjɔn reˈkaz ɛz ɪn kɔjˈmaθ lɔnk}{{\Def=\Dist.\Sg} {\First\Sg.\Agt} {require.\Pst.\Ptcp} be {\Indf} holiday long}{What I need is a long holiday.}
With a verb whose experiencer is an oblique, such as \bl{stovir} `be needed', the free relative has a present participle and an oblique clitic instead: \bl{l'oç me stovent} `what I need'.

Free-choice `whoever' and `whatever' are \bl{cuscon} and \bl{calscon} (Section~\ref{sec:indefinite-pronouns}). \bl{Cuscon} takes a present participle, like a subject relative, and \bl{calscon} an agentive pronoun and a past participle, like an object relative. As neither contains a finite verb, there is no choice of mood to be made.
    \glsen{Cuscon lo trovant porra lo vardir.}{kiˈxɔn lo troˈvant pɔˈra lo varˈdɪr}{whoever {\Third\Sg.\N.\Acc} {find-\Prs.\Ptcp} {can-\Fut} {\Third\Sg.\N.\Acc} {keep-\Inf}}{Whoever finds it can keep it.}
    \glsen{Calscon teyon deciduð, jo vil t'aïðar.}{kalˈxɔn tiˈjɔn ˌde.tsiˈdɪθ ʒo vɪl ˌtɛj.iˈðar}{whatever {\Second\Sg.\Agt} {decide.\Pst.\Ptcp} {\First\Sg} will {\Second\Sg.\Obl=help-\Inf}}{Whatever you decide, I'll help you.}
A concessive free relative standing before the main clause does not cause inversion (\bl{jo vil}, not \bl{vil jo}).
% Older texts use l'ig for the free relative 'what'; that use is deprecated.

\section{Relativising out of an embedded clause}\label{sec:relative-islands}
A relative clause cannot be formed on a constituent of a clause embedded inside it. English \textit{the man I think you met} has no direct equivalent. Instead the verb of thinking or saying is added as a parenthetical, \bl{crey jo} `I believe' (with the inversion typical of such asides, which is somewhat formal) or \bl{lonc mey} `according to me', and the relative clause is formed on the lower verb alone.
    \glsen{Vey l'hom oç, crey jo, teyon encontrað.}{vi lɔm ɔts kri ʒo tiˈjɔn ˌɛn.kɔnˈtraθ}{look {\Def=person} {\Dist.\Sg} believe {\First\Sg} {\Second\Sg.\Agt} {meet-\Pst.\Ptcp}}{That's the man I think you met.}
The presentational \bl{Vey} `see', the imperative of \bl{veir} (polite \bl{Veyað}), is the usual way of pointing something out, literally or figuratively.
% CHECK: resumptive pronouns are not attested in any relative clause. Is a
% resumptive ever possible, e.g. inside a dont clause or in very informal yon
% relatives?
